\documentclass[11pt,letterpaper,spanish]{article}

\usepackage{caption}

\usepackage{iftex}
\ifPDFTeX
  % --- Setup for pdfLaTeX ---
  \usepackage[utf8]{inputenc}
  \usepackage[T1]{fontenc}
\else
  % --- Setup for XeLaTeX & LuaLaTeX ---
  \usepackage{fontspec}
  \setmainfont{Latin Modern Roman}
\fi


\usepackage[spanish]{babel}

\usepackage{url}
\usepackage{hyperref}
\usepackage{geometry}
\usepackage{svg}
\usepackage{fancyhdr}
\usepackage{xcolor}
\usepackage{multicol}
\usepackage{enumitem}

%% Paquetes caseros
\usepackage{customcode}

% ==========================================
% MARCADO DE PROCEDENCIA DE LAS JUSTIFICACIONES
% (permite distinguir lo tomado de Explicación.pdf
%  de las propuestas de mejora para poder depurarlas)
% ==========================================
\definecolor{propcolor}{RGB}{0,90,160}

\tcbuselibrary{breakable}

\newcommand{\marcaOriginal}{}
\newcommand{\marcaPropuesta}{\textnormal{\small\textcolor{propcolor}{(propuesta)}}}

\newtcolorbox{cajapropuesta}{
	colback=propcolor!4,
	colframe=propcolor,
	boxrule=0.5pt,
	arc=2pt,
	left=5pt, right=5pt, top=5pt, bottom=5pt,
	breakable
}



% Required by biblatex-apa for context-sensitive quotes
\usepackage{csquotes}

% Load biblatex with official APA 7th edition style
\usepackage[style=apa, backend=biber]{biblatex}
\addbibresource{referencias.bib}


\geometry{margin=2cm, top=2.5cm, bottom=2.5cm}

\definecolor{headergray}{gray}{0.3}
\definecolor{rulecolor}{gray}{0.6}

\pagestyle{fancy}
\fancyhf{}
\renewcommand{\headrulewidth}{0.5pt}
\renewcommand{\footrulewidth}{0.5pt}
\fancyhead[L]{\small\textcolor{headergray}{UNIVERSIDAD NACIONAL AUTÓNOMA DE MÉXICO}}
\fancyhead[R]{\small\textcolor{headergray}{\documentoTitulo}}
\fancyfoot[L]{\small\textcolor{headergray}{FACULTAD DE CIENCIAS, 2027-I}}
\fancyfoot[C]{\small\textcolor{headergray}{\thepage}}
\fancyfoot[R]{\small\textcolor{headergray}{FUNDAMENTOS DE BASES DE DATOS}}

% ==========================================
% CONFIGURACIÓN DE VARIABLES (EDITAR AQUÍ)
% ==========================================
\newcommand{\documentoTitulo}{Practica 04}
\newcommand{\documentoSubtitulo}{Modelo Relacional}

\newcommand{\profesorNombre}{Gerardo Avilés Rosas}

\newcommand{\ayudanteTeoriaUno}{Oscar Daniel Flores Linares}
\newcommand{\ayudanteTeoriaDos}{Claudia Itzel Gutiérrez Sánchez}

\newcommand{\ayudanteLabUno}{Ricardo Badillo Macías}
\newcommand{\ayudanteLabDos}{Rocío Aylin Huerta González}

\newcommand{\equipoNombre}{Dobles Comillas}

\newcommand{\integranteUno}{Contreras Morales Román Ismael}
\newcommand{\integranteDos}{Cornejo de la mora Iñaki}
\newcommand{\integranteTres}{Hipólito Felipe Christopher Guillermo}
\newcommand{\integranteCuatro}{Ledesma Cuevas Eduardo}

% ==========================================

\begin{document}

\thispagestyle{empty}

\begin{center}
	\vspace*{1cm}

	\begin{tabular}{cc}
		\includegraphics[width=3cm]{escudos/unam_seal.pdf} & \includegraphics[width=3cm]{escudos/fc_seal.pdf}
	\end{tabular}

	\vspace{0.8cm}
	{\LARGE\bfseries UNIVERSIDAD NACIONAL AUTÓNOMA DE MÉXICO}\\[0.3cm]
	{\large FACULTAD DE CIENCIAS, 2027-I}\\[0.3cm]
	{\large\bfseries FUNDAMENTOS DE BASES DE DATOS}

	\vspace{1.2cm}
	{\huge\bfseries \documentoTitulo:}\\[0.3cm]
	{\Large \documentoSubtitulo}

	\vspace{1.5cm}
\end{center}

\begin{multicols}{2}
	\raggedcolumns

	\textbf{PROFESOR:}\\
	\profesorNombre\\[1.2cm]

	\textbf{AYUDANTES DE TEORÍA:}\\
	\ayudanteTeoriaUno\\
	\ayudanteTeoriaDos\\[1.2cm]

	\textbf{AYUDANTES DE LABORATORIO:}\\
	\ayudanteLabUno\\
	\ayudanteLabDos\\[1.2cm]

	\columnbreak

	\textbf{EQUIPO:}\\
	\equipoNombre\\[1.2cm]

	\textbf{INTEGRANTES DEL EQUIPO:}\\
	\integranteUno\\
	\integranteDos\\
	\integranteTres\\
	\integranteCuatro

\end{multicols}

\vspace{1cm}

\section{Dominio de atributos y restricciones de cada relación}

\subsection{Cliente}
Su llave primaria es \textbf{\textit{idCliente}}, con dominio \textbf{\textit{int}} y restricción de integridad de entidad. Los atributos \textbf{\textit{Nombre}}, \textbf{\textit{Paterno}} y \textbf{\textit{Materno}} se derivan de desglosar el atributo compuesto original para garantizar que cada dominio sea una unidad indivisible; tienen un dominio \textbf{\textit{varchar(100)}} con restricción \textbf{\textit{NOT NULL}} (excepto \textbf{\textit{Materno}}). El atributo \textbf{\textit{fechaNacimiento}} pertenece al dominio \textbf{\textit{date}} (\textbf{\textit{NOT NULL}}), y \textbf{\textit{Sexo}} es \textbf{\textit{varchar(50)}}.

\subsection{Telefono\_Cliente}
Como viene de un atributo multivaluado, tiene una llave primaria compuesta formada por \textbf{\textit{idCliente}} y \textbf{\textit{telefono}}. El atributo \textbf{\textit{idCliente}} (\textbf{\textit{int}}) actúa además como llave foránea (\textbf{\textit{FK}}), referenciando a la relación Cliente, y está sujeto a restricción de integridad referencial. El atributo \textbf{\textit{telefono}} tiene un dominio \textbf{\textit{varchar(15)}} y restricción \textbf{\textit{NOT NULL}} de dominio.

\subsection{Correo\_Cliente}
Como viene de un atributo multivaluado, su llave primaria es compuesta, integrando el identificador del cliente y el \textbf{\textit{correo}}. La columna \textbf{\textit{idCliente}} (\textbf{\textit{int}}) es una llave foránea hacia Cliente con restricción de integridad referencial. El atributo \textbf{\textit{correo}} (\textbf{\textit{varchar(100)}}) tiene restricción de dominio (\textbf{\textit{NOT NULL}}).

\subsection{Tarjeta}
La relación Tarjeta resulta de la conversión de una entidad débil, por lo que su llave primaria es compuesta, formada por el atributo de la entidad fuerte (\textbf{\textit{idCliente}}) y su propio discriminante (\textbf{\textit{idTarjeta}}). El atributo \textbf{\textit{idCliente}} (\textbf{\textit{int}}) funge también como llave foránea (\textbf{\textit{FK}}) hacia la relación Cliente. El atributo \textbf{\textit{idTarjeta}} pertenece al dominio \textbf{\textit{int}} con restricción \textbf{\textit{NOT NULL}}. El atributo \textbf{\textit{Saldo}} se define como un valor numérico exacto (\textbf{\textit{decimal(8,2)}}) con restricción \textbf{\textit{NOT NULL}}. Finalmente, al tener una relación 1:N (EMITIR) con el cajero, la tabla absorbe como llave foránea el atributo \textbf{\textit{rfc\_cajero}} (\textbf{\textit{varchar(100)}}), garantizando mediante integridad referencial que toda tarjeta esté asociada al empleado que la procesó.

\subsection{Recarga}
Su llave primaria es \textbf{\textit{idRecarga}}, perteneciente al dominio de los enteros (\textbf{\textit{int}}) con restricción de integridad de entidad (\textbf{\textit{NOT NULL}}, \textbf{\textit{UNIQUE}}). Sus atributos \textbf{\textit{MontoAbonar}} (\textbf{\textit{decimal(8,2)}}), \textbf{\textit{fecha}} (\textbf{\textit{date}}), \textbf{\textit{hora}} (\textbf{\textit{time}}) y \textbf{\textit{MetodoPago}} (\textbf{\textit{varchar(50)}}) operan bajo restricciones de integridad de dominio (\textbf{\textit{NOT NULL}}). Estructuralmente, al ser una operación que vincula el lugar, el medio de pago y el empleado que la cobra, la tabla absorbe múltiples llaves foráneas (\textbf{\textit{FK}}): \textbf{\textit{idSucursal}} (\textbf{\textit{int}}) con integridad referencial hacia Sucursal; la llave compuesta de la tarjeta recargada, conformada por \textbf{\textit{idCliente}} (\textbf{\textit{int}}) y \textbf{\textit{idTarjeta}} (\textbf{\textit{int}}); y finalmente \textbf{\textit{rfc\_cajero}} (\textbf{\textit{varchar(100)}}) con integridad referencial hacia la relación Cajero.

\subsection{Sucursal}
Su llave primaria es \textbf{\textit{idSucursal}}, la cual pertenece al dominio de los enteros (\textbf{\textit{int}}) y aplica una restricción de integridad de entidad (\textbf{\textit{NOT NULL}}, \textbf{\textit{UNIQUE}}). El atributo compuesto original de la dirección fue descompuesto en las columnas indivisibles \textbf{\textit{Calle}}, \textbf{\textit{NumExterior}}, \textbf{\textit{NumInterior}}, \textbf{\textit{Colonia}} y \textbf{\textit{Estado}}, todas con dominio \textbf{\textit{varchar(50)}}. A estas columnas de ubicación se suman los atributos descriptivos \textbf{\textit{Nombre}}, \textbf{\textit{Telefono}} y \textbf{\textit{Horario}} (\textbf{\textit{varchar(100)}}); el conjunto cuenta con restricción de integridad de dominio (\textbf{\textit{NOT NULL}}).

\subsection{Maquina}
Su llave primaria es \textbf{\textit{idMaquina}}, perteneciente al dominio de los enteros (\textbf{\textit{int}}) con restricción de integridad de entidad (\textbf{\textit{NOT NULL}}, \textbf{\textit{UNIQUE}}). Como atributo solo conserva \textbf{\textit{Estado}}, con dominio \textbf{\textit{varchar(50)}} y restricción \textbf{\textit{NOT NULL}} de dominio. Por su participación en el lado "N" de dos relaciones distintas (PERTENECE con Sucursal y CONTENER con Juego) y por regla de cardinalidad, la tabla absorbe las llaves de ambas entidades. Por lo tanto, contiene \textbf{\textit{idSucursal}} (\textbf{\textit{int}}) e \textbf{\textit{idJuego}} (\textbf{\textit{int}}) como llaves foráneas (\textbf{\textit{FK}}). Ambas cuentan con restricciones estrictas de integridad referencial (\textbf{\textit{NOT NULL}}).

\subsection{Juego}
Su llave primaria es \textbf{\textit{idJuego}}, la cual pertenece al dominio de los enteros (\textbf{\textit{int}}) con restricción de integridad de entidad (\textbf{\textit{NOT NULL}}, \textbf{\textit{UNIQUE}}). Sus atributos \textbf{\textit{Nombre}} (\textbf{\textit{varchar(100)}}) y \textbf{\textit{Tipo}} (\textbf{\textit{varchar(50)}}) pertenecen al dominio de cadena de caracteres, mientras que el atributo \textbf{\textit{Costo}} pertenece al dominio \textbf{\textit{decimal(8,2)}}; todos ellos cuentan con restricción de integridad de dominio (\textbf{\textit{NOT NULL}}).

\subsection{Partida}
Su llave primaria es \textbf{\textit{idPartida}}, perteneciente al dominio de los enteros (\textbf{\textit{int}}) con restricción de integridad de entidad (\textbf{\textit{NOT NULL}}, \textbf{\textit{UNIQUE}}). Sus atributos \textbf{\textit{Puntaje}} y \textbf{\textit{PuntosGanados}} (ambos en el dominio \textbf{\textit{int}}), así como \textbf{\textit{fecha}} (\textbf{\textit{date}}) y \textbf{\textit{hora}} (\textbf{\textit{time}}), todos protegidos con restricciones de integridad de dominio (\textbf{\textit{NOT NULL}}). Al participar en el lado "N" de dos relaciones distintas (OCURRIR con Maquina y REGISTRAR con Tarjeta), absorbe las llaves primarias de dichas entidades. Por lo tanto, contiene como llaves foráneas (\textbf{\textit{FK}}) a \textbf{\textit{idMaquina}} (\textbf{\textit{int}}) y a la llave compuesta de la tarjeta, conformada por \textbf{\textit{idCliente}} (\textbf{\textit{int}}) y \textbf{\textit{idTarjeta}} (\textbf{\textit{int}}). Estas foráneas mantienen una restricción de integridad referencial.

\subsection{Canje}
Su llave primaria es \textbf{\textit{idCanje}}, del dominio de los enteros (\textbf{\textit{int}}), sujeta a la restricción de integridad de entidad (\textbf{\textit{NOT NULL}}, \textbf{\textit{UNIQUE}}). Sus atributos \textbf{\textit{fecha}} (\textbf{\textit{date}}), \textbf{\textit{hora}} (\textbf{\textit{time}}), \textbf{\textit{PuntosUtilizados}} (\textbf{\textit{int}}) y \textbf{\textit{CantidadPremiosObtenidos}} (\textbf{\textit{int}}) operan bajo restricciones de integridad de dominio estrictas (\textbf{\textit{NOT NULL}}). Participa como el lado "N" (muchos) en: OCURRIR con Sucursal, SOLICITAR con Cliente, INCLUIR con Premio y PROCESAR con EncargadoDePremio. Aplicando la regla de cardinalidad, la tabla absorbe cuatro llaves foráneas (\textbf{\textit{FK}}): \textbf{\textit{idSucursal}} (\textbf{\textit{int}}), \textbf{\textit{idCliente}} (\textbf{\textit{int}}), \textbf{\textit{idPremio}} (\textbf{\textit{int}}) y \textbf{\textit{rfc\_EncargadoDePremios}} (\textbf{\textit{varchar(100)}}). Todas estas foráneas aplican reglas de integridad referencial para asegurar que un canje no pueda existir sin estar vinculado a la sucursal donde ocurrió, el cliente que lo solicitó, el premio entregado y el empleado que lo autorizó.

\subsection{Premio}
Su llave primaria es \textbf{\textit{idPremio}}, perteneciente al dominio de los enteros (\textbf{\textit{int}}), con restricción de integridad de entidad (\textbf{\textit{NOT NULL}}, \textbf{\textit{UNIQUE}}). Sus atributos \textbf{\textit{Nombre}}, \textbf{\textit{RangoEdad}} y \textbf{\textit{Categoria}} utilizan el dominio de cadena de caracteres (\textbf{\textit{varchar(50)}}); \textbf{\textit{PuntosNecesarios}} opera en el dominio de los enteros (\textbf{\textit{int}}); y \textbf{\textit{ValorAprox}} pertenece al dominio numérico exacto (\textbf{\textit{decimal(8,2)}}). Todos estos atributos cuentan con restricciones de integridad de dominio (\textbf{\textit{NOT NULL}}).

\subsection{Inventario}
La relación Inventario resulta de la conversión de una entidad débil, por lo que su llave primaria es compuesta, estructurada por la llave de su entidad fuerte dominante (\textbf{\textit{idSucursal}}) y su atributo discriminante propio (\textbf{\textit{idArticulo}}). El atributo \textbf{\textit{idSucursal}} pertenece al dominio de los enteros (\textbf{\textit{int}}) y actúa simultáneamente como llave foránea (\textbf{\textit{FK}}) hacia la relación Sucursal con integridad referencial. El discriminante \textbf{\textit{idArticulo}} (\textbf{\textit{int}}) aplica restricción de integridad de entidad para diferenciar grupos de artículos dentro de la misma sucursal. Adicionalmente, debido a su relación de cardinalidad 1:N con el catálogo principal, la tabla absorbe a \textbf{\textit{idPremio}} (\textbf{\textit{int}}) como una segunda llave foránea (\textbf{\textit{FK}}). Finalmente, el atributo descriptivo \textbf{\textit{CantidadDisponible}} pertenece al dominio numérico entero (\textbf{\textit{int}}) y opera bajo restricción de integridad de dominio (\textbf{\textit{NOT NULL}}).

\subsection{Mantenimiento}
Como es una entidad débil, su llave primaria es compuesta, estructurada por la llave de su entidad fuerte dominante (\textbf{\textit{idMaquina}}) y su atributo discriminante propio (\textbf{\textit{idMantenimiento}}). El atributo \textbf{\textit{idMaquina}} pertenece al dominio de los enteros (\textbf{\textit{int}}) y actúa simultáneamente como llave foránea (\textbf{\textit{FK}}) hacia la relación Maquina, aplicando una restricción de integridad referencial. El discriminante \textbf{\textit{idMantenimiento}} (\textbf{\textit{int}}) aplica restricción de integridad de entidad para diferenciar los distintos servicios aplicados a un mismo equipo. Sus atributos \textbf{\textit{fecha}} (\textbf{\textit{date}}), \textbf{\textit{Tipo}} (\textbf{\textit{varchar(50)}}) y \textbf{\textit{DescripcionProblema}} (\textbf{\textit{varchar(200)}}) operan bajo restricciones de integridad de dominio (\textbf{\textit{NOT NULL}}). Finalmente, al relacionarse en cardinalidad 1:N con Tecnico, la tabla absorbe a \textbf{\textit{rfc\_Tecnico}} (\textbf{\textit{varchar(100)}}) como una segunda llave foránea (\textbf{\textit{FK}}).

\subsection{Empleado}
La jerarquía de la entidad Empleado presenta una especialización total y disyuntiva; sin embargo, la existencia de múltiples atributos multivaluados (teléfonos y correos en la superclase, y certificaciones en la subclase Tecnico) hace inviable la eliminación de la tabla Empleado, pues si la eliminamos, tendremos 2 tablas más (las de atributos multivaluados) por cada subentidad.

Se genera la relación Empleado que concentra los atributos atómicos comunes, estableciendo \textbf{\textit{rfc}} (\textbf{\textit{varchar(100)}}) como su llave primaria (\textbf{\textit{PK}}). Sus atributos son \textbf{\textit{CURP}}, \textbf{\textit{Nombre}}, \textbf{\textit{Paterno}} y \textbf{\textit{Materno}} (\textbf{\textit{varchar(100)}}), \textbf{\textit{FechaNacimiento}} y \textbf{\textit{FechaContratacion}} (\textbf{\textit{date}}), y \textbf{\textit{Salario}} (\textbf{\textit{decimal(8,2)}}), todos bajo restricción de integridad de dominio (\textbf{\textit{NOT NULL}}). Los atributos multivaluados de la superclase se extraen en las relaciones débiles \textbf{\textit{Telefono\_Empleado}} y \textbf{\textit{Correo\_Empleado}}, estructuradas con llaves primarias compuestas que incluyen una llave foránea (\textbf{\textit{FK}}).

Las subentidades (Gerente, Cajero, EncargadoDePremio y Tecnico) se mapean como relaciones independientes que absorben el \textbf{\textit{rfc}} simultáneamente como llave primaria (\textbf{\textit{PK}}) y llave foránea (\textbf{\textit{FK}}) hacia la superclase, añadiendo únicamente sus atributos atómicos específicos (\textbf{\textit{CedulaProfesional}}, \textbf{\textit{HorarioTrabajo}}, \textbf{\textit{AreaEspecialidad}}). El atributo multivaluado del técnico se aísla en la tabla \textbf{\textit{Certificacion\_Tecnico}} aplicando la misma regla de llave compuesta. Finalmente, respetando las cardinalidades del modelo conceptual, la relación Gerente absorbe a \textbf{\textit{idSucursal}} (\textbf{\textit{int}}) como llave foránea (\textbf{\textit{FK}}) bajo restricción \textbf{\textit{UNIQUE}}, por su participación total en la relación 1:1 Administrar con Sucursal.

\subsection{Telefono\_Empleado}
Como viene de un atributo multivaluado de la superclase Empleado, su llave primaria es compuesta por \textbf{\textit{rfc}} (\textbf{\textit{varchar(100)}}) y \textbf{\textit{telefono}} (\textbf{\textit{varchar(15)}}). La columna \textbf{\textit{rfc}} es una llave foránea (\textbf{\textit{FK}}) hacia la relación Empleado con restricción de integridad referencial, y \textbf{\textit{telefono}} tiene restricción \textbf{\textit{NOT NULL}} de dominio.

\subsection{Correo\_Empleado}
Análoga a Telefono\_Empleado, proviene de un atributo multivaluado de la superclase; su llave primaria es compuesta por \textbf{\textit{rfc}} (\textbf{\textit{varchar(100)}}) y \textbf{\textit{correo}} (\textbf{\textit{varchar(100)}}). La columna \textbf{\textit{rfc}} es una llave foránea (\textbf{\textit{FK}}) hacia Empleado con restricción de integridad referencial, y \textbf{\textit{correo}} tiene restricción \textbf{\textit{NOT NULL}} de dominio.

\subsection{Cajero}
Subentidad de Empleado. Su llave primaria es \textbf{\textit{rfc}} (\textbf{\textit{varchar(100)}}), que actúa simultáneamente como llave foránea (\textbf{\textit{FK}}) hacia Empleado con integridad referencial. Añade el atributo propio \textbf{\textit{HorarioTrabajo}} (\textbf{\textit{varchar(100)}}, \textbf{\textit{NOT NULL}}). Participa en el lado "1" de las relaciones 1:N EMITIR (con Tarjeta) y PROCESAR (con Recarga).

\subsection{Gerente}
Subentidad de Empleado. Su llave primaria es \textbf{\textit{rfc}} (\textbf{\textit{varchar(100)}}), que actúa además como llave foránea (\textbf{\textit{FK}}) hacia Empleado. Añade el atributo propio \textbf{\textit{CedulaProfesional}} (\textbf{\textit{varchar(100)}}, \textbf{\textit{NOT NULL}}). Por su participación total en la relación 1:1 ADMINISTRAR con Sucursal, absorbe a \textbf{\textit{idSucursal}} (\textbf{\textit{int}}) como llave foránea (\textbf{\textit{FK}}) bajo restricción \textbf{\textit{UNIQUE}}, con integridad referencial hacia Sucursal.

\subsection{EncargadoDePremio}
Subentidad de Empleado. Su llave primaria es \textbf{\textit{rfc}} (\textbf{\textit{varchar(100)}}), que actúa además como llave foránea (\textbf{\textit{FK}}) hacia Empleado. Añade el atributo propio \textbf{\textit{HorarioTrabajo}} (\textbf{\textit{varchar(100)}}, \textbf{\textit{NOT NULL}}). Participa en el lado "1" de la relación 1:N PROCESAR con Canje.

\subsection{Tecnico}
Subentidad de Empleado. Su llave primaria es \textbf{\textit{rfc}} (\textbf{\textit{varchar(100)}}), que actúa además como llave foránea (\textbf{\textit{FK}}) hacia Empleado. Añade el atributo propio \textbf{\textit{AreaEspecialidad}} (\textbf{\textit{varchar(100)}}, \textbf{\textit{NOT NULL}}). Participa en el lado "1" de la relación 1:N REALIZAR con Mantenimiento, y su atributo multivaluado de certificaciones se aísla en la relación \textbf{\textit{Certificacion\_Tecnico}}.

\subsection{Certificacion\_Tecnico}
Proviene del atributo multivaluado Certificaciones de la subentidad Tecnico; su llave primaria es compuesta por \textbf{\textit{rfcTecnico}} (\textbf{\textit{varchar(100)}}) y \textbf{\textit{certificacion}} (\textbf{\textit{varchar(100)}}). La columna \textbf{\textit{rfcTecnico}} es una llave foránea (\textbf{\textit{FK}}) hacia la relación Tecnico con restricción de integridad referencial, y \textbf{\textit{certificacion}} tiene restricción \textbf{\textit{NOT NULL}} de dominio.



% ==========================================
% Diagrama (página de mayor tamaño)
% ==========================================
\clearpage
% Página de mayor tamaño: se fija el tamaño físico de la página y se centra
% un bloque de texto ligeramente menor (1 cm de margen por lado).
\pdfpagewidth=35cm
\pdfpageheight=35cm
\setlength{\textwidth}{34cm}
\setlength{\hsize}{34cm}
\setlength{\oddsidemargin}{\dimexpr(\pdfpagewidth-\textwidth)/2 - 1in\relax}
\setlength{\evensidemargin}{\oddsidemargin}
\setlength{\textheight}{34cm}
\setlength{\vsize}{34cm}
\setlength{\headheight}{0pt}
\setlength{\headsep}{0pt}
\setlength{\topmargin}{\dimexpr(\pdfpageheight-\textheight)/2 - 1in\relax}
\thispagestyle{empty}
\subsection*{Diagrama Modelo Relacional}
\vspace*{\fill}
\centering
\includegraphics[width=\textwidth, height=0.90\textheight, keepaspectratio]{figuras/relacional.pdf}
\captionof{figure}{Modelo Relacional del centro de entretenimiento familiar PuellaGame.}
\label{fig:modelo_relacional}
\vspace*{\fill}
\clearpage
\restoregeometry
\pdfpagewidth=\paperwidth
\pdfpageheight=\paperheight

\printbibliography

\end{document}
